\chapter{خوشه‌بندی}
\label{ch:clustering}

\section*{اهداف این فصل}
\begin{itemize}
    \item درک مفهوم خوشه‌بندی و تفاوت آن با طبقه‌بندی
    \item یادگیری الگوریتم k-means و نحوهٔ انتخاب مقدار $k$
    \item آشنایی با خوشه‌بندی سلسله‌مراتبی و روش‌های پیوند
    \item آشنایی با DBSCAN برای خوشه‌های با شکل دلخواه
    \item آشنایی با معیارهای ارزیابی خوشه‌بندی
\end{itemize}

\section{مقدمه‌ای بر خوشه‌بندی}

\dmterm{\textbf{خوشه‌بندی}}{Clustering} دومین وظیفهٔ اصلی یادگیری بدون
نظارت است. هدف آن، گروه‌بندی نمونه‌ها به‌گونه‌ای است که نمونه‌های درون
یک گروه (\dmterm{خوشه}{Cluster}) شباهت زیادی به هم داشته باشند و
نمونه‌های دو گروه مختلف، شباهت کمی به یکدیگر داشته باشند — بدون آنکه
هیچ برچسب از پیش‌تعیین‌شده‌ای در اختیار باشد. این وظیفه با طبقه‌بندی
تفاوت اساسی دارد: در طبقه‌بندی، کلاس‌ها از پیش مشخص‌اند و مدل یاد
می‌گیرد نمونهٔ جدید به کدام کلاس تعلق دارد؛ در خوشه‌بندی، حتی تعداد و
ماهیت گروه‌ها نیز باید از دل خود داده کشف شود.

کاربردهای رایج خوشه‌بندی عبارت‌اند از بخش‌بندی مشتریان بر اساس رفتار
خرید، فشرده‌سازی تصویر، و به‌عنوان گامی مقدماتی پیش از تحلیل‌های دیگر
(مثلاً یافتن زیرگروه‌های همگن در یک جمعیت پیش از مدل‌سازی جداگانه برای
هر زیرگروه).

\section{الگوریتم k-means}

\dmterm{\textbf{الگوریتم} \LR{k-means}}{k-means Algorithm} پرکاربردترین
الگوریتم خوشه‌بندی است. ایدهٔ آن ساده است: داده به $k$ خوشه تقسیم
می‌شود، به‌گونه‌ای که هر نمونه به نزدیک‌ترین
\dmterm{مرکز خوشه}{Centroid} تعلق داشته باشد. الگوریتم به‌صورت
تکرارشونده عمل می‌کند:

\begin{enumerate}
    \item ابتدا $k$ مرکز خوشه به‌طور تصادفی انتخاب می‌شوند.
    \item هر نمونه به نزدیک‌ترین مرکز خوشه (بر اساس فاصلهٔ اقلیدسی)
          نسبت داده می‌شود.
    \item مرکز هر خوشه، به‌عنوان میانگین همهٔ نمونه‌های آن خوشه، دوباره
          محاسبه می‌شود.
    \item مراحل ۲ و ۳ تا زمانی تکرار می‌شوند که مراکز خوشه دیگر تغییر
          معناداری نکنند (همگرایی).
\end{enumerate}

هدف k-means کمینه‌کردن مجموع مربعات فاصلهٔ درون‌خوشه‌ای است
\footnote{\LR{Within-Cluster Sum of Squares (WCSS)}}. نکات مهم عملی:

\begin{itemize}
    \item چون انتخاب اولیهٔ مراکز تصادفی است، نتیجهٔ نهایی می‌تواند به
          این انتخاب حساس باشد؛ در عمل الگوریتم چند بار با مقداردهی
          اولیهٔ متفاوت اجرا و بهترین نتیجه انتخاب می‌شود.
    \item k-means برای خوشه‌های با شکل تقریباً کروی و اندازهٔ مشابه بهتر
          عمل می‌کند و برای خوشه‌های با شکل دلخواه یا اندازهٔ بسیار
          متفاوت مناسب نیست (بخش \ref{sec:dbscan} را ببینید).
    \item از آنجا که بر پایهٔ فاصله عمل می‌کند، نرمال‌سازی ویژگی‌ها
          (فصل \ref{ch:data}) پیش از اجرا ضروری است.
\end{itemize}

\subsection{انتخاب مقدار k}

برخلاف طبقه‌بندی که تعداد کلاس‌ها از قبل مشخص است، در k-means باید
مقدار $k$ را از قبل تعیین کنیم. دو روش رایج برای این کار عبارت‌اند از:

\paragraph{روش آرنج}
در \dmterm{\textbf{روش آرنج}}{Elbow Method}، الگوریتم برای مقادیر
مختلف $k$ اجرا می‌شود و مقدار \LR{WCSS} برای هر $k$ رسم می‌شود. با
افزایش $k$، این مقدار همواره کاهش می‌یابد، اما نرخ کاهش معمولاً در یک
نقطهٔ خاص به‌شدت کم می‌شود و نموداری شبیه آرنج بازو ایجاد می‌کند؛ آن
نقطه، انتخاب مناسبی برای $k$ است.

\paragraph{ضریب سیلوئت}
\dmterm{\textbf{ضریب سیلوئت}}{Silhouette Coefficient} برای هر نمونه،
میزان شباهت آن به خوشهٔ خودش را در مقایسه با نزدیک‌ترین خوشهٔ دیگر
می‌سنجد و مقداری بین $-1$ تا $1$ می‌دهد؛ مقدار نزدیک به $1$ نشان‌دهندهٔ
خوشه‌بندی مناسب است. میانگین این ضریب برای همهٔ نمونه‌ها، معیاری کلی
برای مقایسهٔ کیفیت خوشه‌بندی با مقادیر مختلف $k$ به دست می‌دهد و
برخلاف روش آرنج، نیازی به قضاوت چشمی ندارد.

\section{خوشه‌بندی سلسله‌مراتبی}

برخلاف k-means که تعداد خوشه‌ها را از پیش می‌خواهد،
\dmterm{\textbf{خوشه‌بندی سلسله‌مراتبی}}{Hierarchical Clustering}
ساختاری درختی از خوشه‌ها در سطوح مختلف می‌سازد که با نموداری به نام
\dmterm{\textbf{دندروگرام}}{Dendrogram} نمایش داده می‌شود؛ تعداد نهایی
خوشه‌ها را می‌توان بعداً با «بریدن» این درخت در ارتفاع دلخواه تعیین
کرد. رایج‌ترین رویکرد،
\dmterm{\textbf{رویکرد تجمیعی}}{Agglomerative Approach} است که به‌صورت
پایین به بالا عمل می‌کند: در ابتدا هر نمونه یک خوشهٔ جداگانه است، و در
هر مرحله، دو خوشهٔ نزدیک‌تر به هم ادغام می‌شوند تا در نهایت یک خوشهٔ
واحد باقی بماند.

\subsection{روش‌های پیوند}

معیار «نزدیکی» میان دو خوشه (نه دو نمونهٔ منفرد) با
\dmterm{\textbf{روش‌های پیوند}}{Linkage Methods} تعریف می‌شود:

\begin{itemize}
    \item \dmterm{\textbf{پیوند تکی}}{Single Linkage}: کمترین فاصله میان
          هر عضو یک خوشه با هر عضو خوشهٔ دیگر.
    \item \dmterm{\textbf{پیوند کامل}}{Complete Linkage}: بیشترین فاصله
          میان اعضای دو خوشه.
    \item \dmterm{\textbf{پیوند میانگین}}{Average Linkage}: میانگین
          فاصلهٔ همهٔ جفت‌های ممکن میان اعضای دو خوشه.
\end{itemize}

پیوند تکی می‌تواند به پدیدهٔ \dmterm{زنجیره‌ای‌شدن}{Chaining} (تشکیل
خوشه‌های کشیده و نامناسب) منجر شود، در حالی که پیوند کامل معمولاً خوشه‌های
فشرده‌تر و کروی‌تری تولید می‌کند.

\section{DBSCAN}
\label{sec:dbscan}

\dmterm{\textbf{الگوریتم} \LR{DBSCAN}}{Density-Based Spatial Clustering
of Applications with Noise} برخلاف دو روش قبل، خوشه‌ها را بر اساس
\textbf{چگالی} نقاط تعریف می‌کند، نه فاصله تا یک مرکز یا ساختار
سلسله‌مراتبی. این ویژگی به آن اجازه می‌دهد خوشه‌هایی با شکل دلخواه
(نه‌فقط کروی) را شناسایی کند و به‌طور طبیعی، نقاطی که در هیچ ناحیهٔ
پرچگالی قرار نمی‌گیرند را به‌عنوان \textbf{نویز} یا داده پرت علامت‌گذاری
کند (که ارتباط مستقیمی با فصل \ref{ch:outliers} دارد).

DBSCAN دو پارامتر اصلی دارد: شعاع همسایگی ($\varepsilon$) و حداقل تعداد
نقاط لازم در آن همسایگی برای تشکیل یک ناحیهٔ پرچگالی\footnote{\LR{MinPts}}.
نقاطی که این حداقل تعداد همسایه را در شعاع خود
داشته باشند \dmterm{نقاط هسته‌ای}{Core Points} نامیده می‌شوند و خوشه‌ها
با اتصال نقاط هسته‌ای مجاور به یکدیگر شکل می‌گیرند.

\section{ارزیابی خوشه‌بندی}

از آنجا که در خوشه‌بندی برچسب واقعی وجود ندارد، ارزیابی آن دشوارتر از
طبقه‌بندی است. رایج‌ترین رویکرد، استفاده از همان
\dmterm{ضریب سیلوئت}{Silhouette Coefficient} است که در بخش انتخاب $k$
دیدیم؛ این معیار می‌تواند برای مقایسهٔ کیفیت هر خوشه‌بندی (فارغ از
الگوریتم تولیدکننده) به کار رود. در مواردی که برچسب واقعی (برای اهداف
پژوهشی) در دسترس باشد، معیارهایی مانند
\dmterm{\textbf{شاخص رند تعدیل‌شده}}{Adjusted Rand Index} نیز برای
مقایسهٔ خوشه‌بندی به‌دست‌آمده با برچسب واقعی استفاده می‌شوند.

\section*{پیاده‌سازی در پایتون}

پیاده‌سازی \texttt{KMeans}، \texttt{AgglomerativeClustering} و
\texttt{DBSCAN} از \texttt{scikit-learn}، به‌همراه محاسبهٔ ضریب سیلوئت
و رسم دندروگرام، در بخش «فصل ۷» پیوست الف (\ref{app:python}) آمده است.

\section*{خلاصهٔ فصل}

در این فصل با خوشه‌بندی به‌عنوان دومین وظیفهٔ اصلی یادگیری بدون نظارت
آشنا شدیم. سه رویکرد اصلی را بررسی کردیم: k-means (مبتنی بر فاصله تا
مرکز)، خوشه‌بندی سلسله‌مراتبی (مبتنی بر ساختار درختی و روش‌های پیوند)
و DBSCAN (مبتنی بر چگالی). همچنین با روش آرنج و ضریب سیلوئت برای انتخاب
تعداد خوشه‌ها و ارزیابی کیفیت خوشه‌بندی آشنا شدیم. در فصل بعد، به
موضوع نزدیک به خوشه‌بندی، یعنی تشخیص داده‌های پرت، می‌پردازیم.

\section*{تمرین‌ها}

\begin{enumerate}
    \item تفاوت اصلی خوشه‌بندی با طبقه‌بندی را از نظر نوع داده مورد نیاز
          (برچسب‌دار در برابر بدون برچسب) توضیح دهید.
    \item چرا نتیجهٔ k-means می‌تواند در دو اجرای مختلف با یک مجموعه
          داده ثابت، متفاوت باشد؟ چگونه می‌توان این مشکل را کاهش داد؟
    \item تفاوت پیوند تکی و پیوند کامل را با یک مثال از شکل خوشه‌های
          حاصل توضیح دهید.
    \item چرا DBSCAN می‌تواند خوشه‌هایی با شکل دلخواه را شناسایی کند در
          حالی که k-means معمولاً خوشه‌های کروی تولید می‌کند؟
    \item توضیح دهید ضریب سیلوئت چگونه هم برای انتخاب مقدار $k$ در
          k-means و هم برای مقایسهٔ کلی کیفیت خوشه‌بندی به کار می‌رود.
\end{enumerate}
